% Standalone problem document, generated from the adjacent README.md.
% Compile with XeLaTeX. Mathematical content is maintained in Markdown.
\documentclass[11pt,a4paper]{article}
\usepackage[fontsize=10.5pt]{scrextend}
\usepackage[margin=22mm,headheight=15pt,headsep=9mm,footskip=11mm]{geometry}
\usepackage{fontspec}
\setmainfont{texgyrepagella}[Extension=.otf,UprightFont=*-regular,BoldFont=*-bold,ItalicFont=*-italic,BoldItalicFont=*-bolditalic]
\setsansfont{texgyreheros}[Extension=.otf,UprightFont=*-regular,BoldFont=*-bold,ItalicFont=*-italic,BoldItalicFont=*-bolditalic]
\setmonofont{texgyrecursor}[Extension=.otf,UprightFont=*-regular,BoldFont=*-bold,ItalicFont=*-italic,BoldItalicFont=*-bolditalic,Scale=MatchLowercase]
\usepackage{amsmath,amssymb}
\usepackage{unicode-math}
\setmathfont{texgyrepagella-math.otf}
\usepackage{xcolor,microtype,enumitem,needspace,fancyhdr}
\usepackage{bookmark}
\definecolor{ink}{HTML}{17344B}
\definecolor{accent}{HTML}{197D80}
\definecolor{muted}{HTML}{596773}
\hypersetup{colorlinks=true,linkcolor=accent,urlcolor=accent,pdftitle={MI-14 - The
complex Lu--Wenzel spectral
conjecture},pdfauthor={Open Problems in Numerical Linear Algebra}}
\urlstyle{same}
\setlength{\parindent}{0pt}
\setlength{\parskip}{5pt plus 1pt minus 1pt}
\linespread{1.04}
\setlength{\emergencystretch}{3em}
\widowpenalty=10000
\clubpenalty=10000
\displaywidowpenalty=10000
\allowdisplaybreaks[1]
\setlist{nosep,leftmargin=1.5em}
\providecommand{\tightlist}{\setlength{\itemsep}{0pt}\setlength{\parskip}{0pt}}
\providecommand{\pandocbounded}[1]{#1}
\pagestyle{fancy}
\fancyhf{}
\fancyhead[L]{\sffamily\footnotesize\color{muted}OPEN PROBLEMS IN NUMERICAL LINEAR ALGEBRA}
\fancyhead[R]{\sffamily\bfseries\footnotesize\color{accent}MI-14}
\fancyfoot[L]{\parbox[b]{0.9\textwidth}{\sffamily\fontsize{8}{10}\selectfont\color{muted}Editorial ratings; a bounded search cannot certify that a problem remains unsolved.\\Literature check: 2026-09-08}}
\fancyfoot[R]{\sffamily\footnotesize\color{muted}\thepage}
\renewcommand{\headrulewidth}{0.3pt}
\renewcommand{\headrule}{\hbox to\headwidth{\color{accent}\leaders\hrule height \headrulewidth\hfill}}
\makeatletter
\renewcommand{\section}{\Needspace{5\baselineskip}\@startsection{section}{1}{0pt}{12pt plus 2pt minus 2pt}{4pt}{\sffamily\bfseries\large\color{ink}}}
\renewcommand{\subsection}{\Needspace{4\baselineskip}\@startsection{subsection}{2}{0pt}{9pt plus 1pt minus 1pt}{3pt}{\sffamily\bfseries\normalsize\color{ink}}}
\makeatother
\setcounter{secnumdepth}{0}
\begin{document}
{\sffamily\small\color{muted}Matrix inequalities and norms\par}
\vspace{3pt}
{\raggedright\hyphenpenalty=10000\exhyphenpenalty=10000\sffamily\bfseries\fontsize{21}{24}\selectfont\color{ink}The
complex Lu--Wenzel spectral conjecture\par}
\vspace{7pt}
{\sffamily\small\textbf{Difficulty:} extreme\quad\textbf{Importance:} interesting
to the community\par}
\vspace{4pt}
{\color{accent}\hrule height 0.8pt}
\vspace{6pt}
\textbf{Provenance:} explicit conjecture

\subsection{Problem statement}\label{problem-statement}

For every integer \(n\ge2\) and \(X\in\mathbb C^{n\times n}\) with
\(\|X\|_F=1\), consider the complex-linear operator

\[T_X:\mathbb C^{n\times n}\longrightarrow\mathbb C^{n\times n},
\qquad T_X(Y)=[X^*,[X,Y]],\qquad [P,Q]=PQ-QP.\]

Use the inner product \(\langle Y,Z\rangle_F=\operatorname{tr}(Y^*Z)\).
The operator is Hermitian positive semidefinite; list its \(n^2\)
eigenvalues, including multiplicities, as
\(\lambda_1(T_X)\ge\cdots\ge\lambda_{n^2}(T_X)\ge0\). Is

\[\sum_{i=1}^{2k}\lambda_i(T_X)\le 2k+2
\qquad\text{for every integer }1\le k\le\lfloor n^2/2\rfloor?\]

\subsection{Relevance}\label{relevance}

This bounds sums of squared singular values of the commutator map
\(Y\mapsto XY-YX\). Such a map is a basic Sylvester operator. The
conjecture strengthens a bound on its largest singular value to
simultaneous bounds on its singular spectrum.

\subsection{References}\label{references}

\begin{enumerate}
\def\labelenumi{\arabic{enumi}.}
\tightlist
\item
  J. Ge, F. Li, Z. Lu, and Y. Zhou, \emph{On some conjectures by Lu and
  Wenzel}, arXiv:1908.06624v1, Conjecture 7; Theorems 1.1--1.2 and §4.
  \href{https://arxiv.org/html/1908.06624v1}{Primary manuscript}.
\item
  J. Ge, F. Li, Z. Tang, and Y. Zhou, \emph{A survey on the DDVV-type
  inequalities}, Advances in Mathematics (China) 53 (2024), 449--467:
  published Conjecture 4.6 and Theorems 4.1--4.2, pp.463--464. These are
  Conjecture 4.14 and Theorems 4.17--4.18 in
  \href{https://arxiv.org/html/2402.01085v1}{arXiv:2402.01085v1}.
  \href{https://ccj.pku.edu.cn/Article/DownLoad?id=374327987&type=ArticleFile}{Published
  PDF}.
\item
  Z. Liu, \emph{The Lu--Wenzel Spectral Conjecture for Square-Zero
  Matrices}, Zenodo record 21693383, abstract.
  \href{https://zenodo.org/records/21693383}{Author's preprint record}.
\end{enumerate}

\subsection{Status check --- 2026-09-08}\label{status-check-2026-09-08}

Both arXiv sources remain v1. The published survey retains the
conjecture and records the normal, rank-one, and \(n=2,3\) cases. Liu's
record claims the additional class \(X^2=0\) and scalar translates; its
abstract does not claim the unrestricted assertion. That record's
displayed publication date and July 2026 creation metadata differ, so no
priority date is inferred from it. Searches used
\textit{Lu-Wenzel conjecture proof 2026},
\textit{Lu Wenzel conjecture solved}, and
\textit{Lu-Wenzel spectral conjecture}. No general resolution was
located. Equivalent majorization and fundamental-commutator
formulations, and the real restriction, are grouped in this entry.
\end{document}
